\documentclass[10pt,twocolumn,letterpaper]{article}
\usepackage{cvpr}
\usepackage{times}
\usepackage{graphicx}
\usepackage{amsmath,amssymb}
\def\paperID{*****}
\def\confName{CVPR}
\def\confYear{2027}
\begin{document}
\onecolumn
\noindent\textbf{ContextWorld overview}

\medskip
\noindent\includegraphics[width=\textwidth]{design/benchmark_overview_pusher_response.pdf}

\smallskip
{\small\noindent\textbf{ContextWorld overview.} (A) A recorded pusher-displacement pair with shared current observations and future actions. Red arrows show action directions; dashed arrows mark observed endpoint displacement. The translucent outline in the lower future panel marks the T block's current pose. These annotations are not model inputs; all scene images use the same crop of the original observations. (B) The frozen model produces latent predictions from each history. The three subpanels illustrate complementary checks for paired tasks. Filled circles denote predictions and open circles denote encoded true futures. The first two views share the same four points: the highlighted point fixes the prediction for future discrimination and the true future for history benefit. Here $d_{ij}$ is the error between the prediction from history $i$ and true future $j$; the inequalities show the desired ordering for both conditions. Response fidelity compares $\Delta\hat z=\hat z_B-\hat z_A$ with $\Delta z=z_B-z_A$, using the same schematic points. All latent-space geometry is illustrative, not a projection of measured model outputs. (C) Nine tasks across four environments.}
\end{document}
